\chapter{Sequências de funções} \label{fs:chapter}

%%%%%%%%%%%%%%%%%%%%%%%%%%%%%%%%%%%%%%%%%%%%%%%%%%%%%%%%%%%%%%%%%%%%%%%%%%%%%%

\section{Convergência pontual e uniforme}
\label{sec:puconv}

%mbxINTROSUBSECTION

\sectionnotes{1--1,5 aulas}

Até agora, quando falamos de limites de sequências, tratamos de sequências
de números. Um conceito muito útil em análise é o de sequência de funções.
Por exemplo, uma solução de certa equação diferencial pode ser encontrada
determinando-se apenas soluções aproximadas. A solução propriamente dita é,
então, algum tipo de limite dessas soluções aproximadas.

Ao tratar de sequências de funções, a dificuldade está no fato de haver
várias noções de limite. Vamos descrever duas noções comuns de limite para
uma sequência de funções.

\subsection{Convergência pontual}

\begin{defn}
\index{convergência pontual}
Para todo $n\in\N$, seja $f_n \colon S \to \R$ uma função. Dizemos que a sequência
$\{ f_n \}_{n=1}^\infty$
\emph{\myindex{converge pontualmente}}\footnote{Salvo indicação em
contrário, \emph{convergir} geralmente significa \emph{convergir
pontualmente}.} para $f \colon S \to \R$ se, para
todo $x\in S$,
\begin{equation*}
f(x) =
\lim_{n\to\infty} f_n(x) .
\end{equation*}
\end{defn}

Os limites de sequências de números são únicos; portanto, se uma sequência
$\{f_n\}_{n=1}^\infty$ converge pontualmente, a função limite $f$ é única.
É comum dizer que $f_n \colon S \to \R$ \emph{converge pontualmente para
$f$ em $T\subset S$}, para alguma função $f \colon T \to \R$. Nesse caso,
queremos dizer que $f(x)=\lim_{n\to\infty}f_n(x)$ para todo $x\in T$. Em
outras palavras, as restrições de $f_n$ a $T$ convergem pontualmente para $f$.


\begin{example}
Em $[-1,1]$, a sequência de funções definida por $f_n(x) \coloneqq x^{2n}$
converge pontualmente para $f \colon [-1,1] \to \R$, dada por
\begin{equation*}
f(x) =
\begin{cases}
1 & \text{se } x=-1 \text{ ou } x=1, \\
0 & \text{caso contrário.}
\end{cases}
\avoidbreak
\end{equation*}
Veja a \figureref{x2nfig}.

\begin{myfigureht}
\myincludegraphics{x2nfig}{%
Gráficos das funções x ao quadrado, x à quarta potência, x à sexta potência
e x à décima sexta potência, no intervalo de menos 1 a 1. Todos começam no
mesmo ponto à esquerda e terminam no mesmo ponto à direita, pois as funções
valem 1 nas extremidades. No interior do intervalo, as potências maiores
deixam o gráfico mais plano e mais próximo do eixo x no centro; perto das
extremidades, porém, ele precisa subir mais rapidamente até 1.}
\caption{Gráficos de $f_1$, $f_2$, $f_3$ e $f_8$, para
$f_n(x) \coloneqq x^{2n}$.\label{x2nfig}}
\end{myfigureht}

Para verificar isso, primeiro suponha que $x\in(-1,1)$. Então
$0\leq x^2<1$. Já vimos que
\begin{equation*}
\sabs{x^{2n} - 0} = {(x^2)}^n \to 0 \quad \text{as} \quad n \to \infty .
\end{equation*}
Portanto, $\lim_{n\to\infty}f_n(x)=0$. Quando $x=1$ ou $x=-1$, temos
$x^{2n}=1$ para todo $n$ e, consequentemente,
$\lim_{n\to\infty}f_n(x)=1$.

Observe que, para $x\notin[-1,1]$, a sequência $\{x^{2n}\}_{n=1}^\infty$
não converge. Assim, se definíssemos $f_n$ em um conjunto maior que
$[-1,1]$, a sequência $\{f_n\}_{n=1}^\infty$ convergiria pontualmente
apenas em $[-1,1]$.
\end{example}

Frequentemente, as funções são dadas por uma série. Nesse caso, usamos a
noção de convergência pontual para determinar os valores da função.

\begin{example} \label{example:geomsumptconv}
Escrevemos
\begin{equation*}
\sum_{k=0}^\infty x^k
\end{equation*}
para denotar o limite das funções
\begin{equation*}
f_n(x) \coloneqq \sum_{k=0}^n x^k .
\end{equation*}
Ao estudar séries, vimos que, em $(-1,1)$, as funções $f_n$ convergem
pontualmente para
\begin{equation*}
\frac{1}{1-x} .
\end{equation*}

O ponto sutil é que, embora $\frac{1}{1-x}$ esteja definida para todo
$x\neq 1$ e $f_n$ estejam definidas para todo $x$ (inclusive $x=1$), a
convergência só ocorre em $(-1,1)$. Portanto, quando escrevemos
\begin{equation*}
f(x) \coloneqq \sum_{k=0}^\infty x^k ,
\end{equation*}
queremos dizer que $f$ está definida em $(-1,1)$ e é o limite pontual das
somas parciais.
\end{example}

\begin{example}
Seja $f_n(x) \coloneqq \sin(nx)$. Então $f_n$ não converge pontualmente
para nenhuma função em qualquer intervalo. A sequência pode convergir em
certos pontos, como $x=0$ ou $x=\pi$. Deixamos como exercício mostrar que,
em qualquer intervalo $[a,b]$, existe um $x$ tal que $\sin(xn)$ não tem
limite quando $n$ tende a infinito. Veja a \figureref{fig:nonconvsinxn}.
\begin{myfigureht}
\myincludegraphics{nonconvsinxn}{%
Gráficos de várias ondas senoidais, com as frequências maiores em cinza mais
claro. É quase como se tentássemos preencher uma faixa horizontal desenhando
linhas onduladas. Todas as curvas começam no mesmo ponto do eixo x, à
esquerda, subindo, e passam pelo mesmo ponto do eixo x à direita, também
subindo. Todas passam pelo centro do eixo x, algumas subindo e outras
descendo. Padrões semelhantes se repetem em outros pontos, mas apenas algumas
das curvas de maior frequência passam por esses pontos.}
\caption{Gráficos de $\sin(nx)$ para $n=1,2,\ldots,10$, com os valores
maiores de $n$ em cinza mais claro.%
\label{fig:nonconvsinxn}}
\end{myfigureht}
\end{example}

Antes de passar à convergência uniforme, vamos reformular a convergência
pontual de outra maneira. Deixamos a demonstração para o leitor: trata-se de
uma aplicação simples da definição de convergência de uma sequência de
números reais.

\begin{prop} \label{ptwsconv:prop}
Sejam $f_n \colon S \to \R$ e $f \colon S \to \R$ funções. Então
$\{f_n\}_{n=1}^\infty$ converge pontualmente para $f$ se, e somente se,
para todo $x\in S$ e todo $\epsilon>0$, existe $N\in\N$ tal que
\begin{equation*}
\babs{f_n(x)-f(x)} < \epsilon \quad \text{para todo } n \geq N .
\end{equation*}
\end{prop}

O ponto essencial é que $N$ pode depender de $x$, e não apenas de
$\epsilon$. Para cada $x$, podemos escolher um $N$ diferente.
Se pudéssemos escolher um único $N$ para todos os valores de $x$, teríamos
o que se chama convergência uniforme.

\subsection{Convergência uniforme}

\begin{defn}
\index{convergência uniforme}
Sejam $f_n \colon S \to \R$ e $f \colon S \to \R$ funções. Dizemos que a
sequência $\{f_n\}_{n=1}^\infty$ \emph{\myindex{converge uniformemente}}
para $f$ se, para todo $\epsilon>0$, existe $N\in\N$ tal que, para todo
$n\geq N$,
\begin{equation*}
\babs{f_n(x) - f(x)} < \epsilon \qquad \text{para todo } x \in S.
\end{equation*}
\end{defn}
\begin{myfigureht}
\myincludegraphics{uniformconv}{%
Gráfico de uma função indicada por f, cercada por uma faixa cujos limites
são indicados por f menos épsilon e f mais épsilon. Uma função indicada por
f índice n aparece como uma linha ondulada inteiramente contida nessa faixa.}
\caption{Na convergência uniforme, para $n\geq N$, as funções $f_n$ ficam
entre os gráficos de $f-\epsilon$ e $f+\epsilon$.%
\label{fig:uniformconv}}
\end{myfigureht}

Na convergência uniforme, $N$ não pode depender de $x$. Dado $\epsilon>0$,
precisamos encontrar um $N$ que sirva para todo $x\in S$. Veja na
\figureref{fig:uniformconv} uma ilustração. A convergência uniforme implica
a convergência pontual, como mostra a demonstração a seguir, que decorre da
\propref{ptwsconv:prop}:

\begin{prop}
Seja $\{f_n\}_{n=1}^\infty$ uma sequência de funções
$f_n \colon S \to \R$. Se $\{f_n\}_{n=1}^\infty$ converge uniformemente
para $f \colon S \to \R$, então $\{f_n\}_{n=1}^\infty$ converge
pontualmente para $f$.
\end{prop}

A recíproca não vale.

\begin{example}
As funções $f_n(x) \coloneqq x^{2n}$ não convergem uniformemente em $[-1,1]$,
embora convirjam pontualmente. Para ver isso, suponha, por contradição, que
a convergência seja uniforme. Para $\epsilon \coloneqq \nicefrac{1}{2}$,
existiria um $N$ tal que $x^{2N}=\sabs{x^{2N}-0}<\nicefrac{1}{2}$ para todo
$x\in(-1,1)$ (pois $f_n(x)$ converge para 0 em $(-1,1)$). Mas isso implica
que, para toda sequência $\{x_k\}_{k=1}^\infty$ em $(-1,1)$ tal que
$\lim_{k\to\infty}x_k=1$, temos $x_k^{2N}<\nicefrac{1}{2}$ para todo $k$.
Por outro lado, $x^{2N}$ é uma função contínua de $x$ (é um polinômio).
Assim, obtemos a contradição
\begin{equation*}
1 = 1^{2N}  = \lim_{k\to\infty} x_k^{2N} \leq \nicefrac{1}{2} .
\end{equation*}

No entanto, se restringirmos o domínio a $[-a,a]$, com $0<a<1$, então
$\{f_n\}_{n=1}^\infty$ converge uniformemente para 0 em $[-a,a]$. Observe
que $a^{2n}\to0$ quando $n\to\infty$. Dado $\epsilon>0$, escolha
$N\in\N$ tal que $a^{2n}<\epsilon$ para todo $n\geq N$.
Se $x\in[-a,a]$, então $\sabs{x}\leq a$. Logo, para todo $n\geq N$ e
todo $x\in[-a,a]$,
\begin{equation*}
\sabs{x^{2n}} = \sabs{x}^{2n} \leq a^{2n} < \epsilon .
\end{equation*}
\end{example}

\subsection{Convergência na norma uniforme}

Para funções limitadas, há outra maneira, mais abstrata, de pensar na
convergência uniforme. A cada função limitada associamos um certo número
não negativo que mede a \myquote{distância} da função à função constante
$0$. Esse número nos permite \myquote{medir} a distância entre duas
funções. Então traduzimos uma afirmação sobre convergência uniforme em uma
afirmação sobre uma certa sequência de números reais que converge a zero.

\begin{defn} \label{def:unifnorm}
Seja $f \colon S \to \R$ uma função limitada. Definimos
\glsadd{not:uniformnorm}
\begin{equation*}
\snorm{f}_S \coloneqq
\sup \Bigl\{ \babs{f(x)} : x \in S \Bigr\} .
\avoidbreak
\end{equation*}
Chamamos $\snorm{\cdot}_S$ de \emph{\myindex{norma uniforme}}.
Às vezes, usa-se outra notação\footnote{A notação e a terminologia não
são totalmente padronizadas. Essa norma também é chamada de
\emph{\myindex{norma do supremo}} ou
\emph{\myindex{norma infinito}}; além de $\snorm{f}_u$ e $\snorm{f}_S$,
às vezes ela é escrita como $\snorm{f}_{\infty}$ ou
$\snorm{f}_{\infty,S}$.}, como $\snorm{f}_u$.
\end{defn}

O subscrito indica o conjunto sobre o qual se toma o supremo.
Assim, se $K \subset S$, então
\begin{equation*}
\snorm{f}_K =
\sup \Bigl\{ \babs{f(x)} : x \in K \Bigr\} .
\end{equation*}

\begin{prop}
Uma sequência de funções limitadas $f_n \colon S \to \R$ converge
uniformemente para $f \colon S \to \R$ se, e somente se,
\begin{equation*}
\lim_{n\to\infty} \snorm{f_n - f}_S = 0 .
\end{equation*}
\end{prop}

\begin{proof}
Suponha primeiro que
$\lim_{n\to\infty} \snorm{f_n - f}_S = 0$. Seja $\epsilon > 0$ dado.
Existe um $N$ tal que, para $n \geq N$, temos
$\snorm{f_n - f}_S < \epsilon$. Como $\snorm{f_n-f}_S$ é o supremo de
$\babs{f_n(x)-f(x)}$, concluímos que, para todo $x \in S$,
$\babs{f_n(x)-f(x)} \leq \snorm{f_n - f}_S < \epsilon$.

Por outro lado, suponha que $\{ f_n \}_{n=1}^\infty$ converge
uniformemente para $f$. Seja $\epsilon > 0$ dado. Então existe $N$ tal
que, para todo $n \geq N$, temos $\babs{f_n(x)-f(x)} < \epsilon$ para
todo $x \in S$. Tomando o supremo em relação a $x \in S$, obtemos
$\snorm{f_n - f}_S \leq \epsilon$. Portanto,
$\lim_{n\to\infty} \snorm{f_n-f}_S = 0$.
\end{proof}

Às vezes, dizemos que
\emph{$\{ f_n \}_{n=1}^\infty$ converge para $f$ na norma uniforme}%
\index{converge na norma uniforme}\index{convergência na norma uniforme}
em vez de dizer que \emph{converge uniformemente}, quando
$\snorm{f_n-f}_S \to 0$. A proposição afirma que, para funções limitadas,
as duas noções são equivalentes.

\begin{example}
Seja $f_n \colon [0,1] \to \R$ dada por
$f_n(x) \coloneqq \frac{nx+ \sin(nx^2)}{n}$.
Vamos mostrar que $\{ f_n \}_{n=1}^\infty$ converge uniformemente para
$f(x) \coloneqq x$. Calculemos:
\begin{equation*}
\begin{split}
\snorm{f_n-f}_{[0,1]}
& =
\sup \left\{ \abs{\frac{nx+ \sin(nx^2)}{n} - x} : x \in [0,1] \right\}
\\
& =
\sup \left\{ \frac{\babs{\sin(nx^2)}}{n} : x \in [0,1] \right\}
\\
& \leq
\sup \bigl\{ \nicefrac{1}{n} : x \in [0,1] \bigr\}
\\
& = \nicefrac{1}{n}.
\end{split}
\end{equation*}
\end{example}

Usando a norma uniforme, definimos sequências de Cauchy de modo semelhante
à definição para sequências de números reais.

\begin{defn}
Sejam $f_n \colon S \to \R$ funções limitadas.
Dizemos que a sequência é \emph{\myindex{de Cauchy na norma uniforme}}
ou \emph{\myindex{uniformemente de Cauchy}}
se, para todo $\epsilon > 0$, existe $N \in \N$ tal que, para todos
$m,k \geq N$,
\begin{equation*}
\snorm{f_m-f_k}_S < \epsilon .
\end{equation*}
\end{defn}

\begin{prop} \label{prop:uniformcauchy}
Sejam $f_n \colon S \to \R$ funções limitadas.
Então $\{ f_n \}_{n=1}^\infty$ é de Cauchy na norma uniforme se, e somente
se existe uma função $f \colon S \to \R$ para a qual
$\{ f_n \}_{n=1}^\infty$ converge uniformemente.
\end{prop}

\begin{proof}
Suponha primeiro que $\{ f_n \}_{n=1}^\infty$ seja de Cauchy na norma
uniforme. Vamos definir $f$. Fixe $x$.
A sequência $\bigl\{ f_n(x) \bigr\}_{n=1}^\infty$ é de Cauchy, pois
\begin{equation*}
\babs{f_m(x)-f_k(x)}
\leq
\snorm{f_m-f_k}_S .
\end{equation*}
Assim, $\bigl\{ f_n(x) \bigr\}_{n=1}^\infty$ converge para algum número real. Defina $f \colon S
\to \R$ por
\begin{equation*}
f(x) \coloneqq \lim_{n \to \infty} f_n(x) .
\end{equation*}
Logo, $\{ f_n \}_{n=1}^\infty$ converge pontualmente para $f$. Para
mostrar que a convergência é uniforme, seja $\epsilon > 0$ dado. Escolha
um $N$ tal que, para todos $m,k \geq N$,
$\snorm{f_m-f_k}_S < \nicefrac{\epsilon}{2}$. Em outras palavras, para
todo $x$, temos $\babs{f_m(x)-f_k(x)} < \nicefrac{\epsilon}{2}$.
Fixado $x$, façamos $k$ tender ao infinito. Então
$\babs{f_m(x)-f_k(x)}$ tende a $\babs{f_m(x)-f(x)}$.
Consequentemente, para todo $x$,
\begin{equation*}
\babs{f_m(x)-f(x)} \leq \nicefrac{\epsilon}{2} < \epsilon .
\end{equation*}
Portanto, $\{ f_n \}_{n=1}^\infty$ converge uniformemente.

Agora, provemos a outra implicação.
Suponha que $\{ f_n \}_{n=1}^\infty$ converge uniformemente para $f$.
Dado $\epsilon > 0$, escolha $N$ tal que, para todo $n \geq N$,
$\babs{f_n(x)-f(x)} < \nicefrac{\epsilon}{4}$ para todo $x \in S$.
Portanto, para todos $m,k \geq N$ e todo $x$,
\begin{multline*}
\babs{f_m(x)-f_k(x)} = 
\babs{f_m(x)-f(x)+f(x)-f_k(x)}
\\
\leq
\babs{f_m(x)-f(x)}+\babs{f(x)-f_k(x)} < \nicefrac{\epsilon}{4} +
\nicefrac{\epsilon}{4} = \nicefrac{\epsilon}{2} .
\end{multline*}
Tomando o supremo em relação a todo $x$, obtemos
\begin{equation*}
\snorm{f_m-f_k}_S \leq \nicefrac{\epsilon}{2} < \epsilon .  \qedhere
\end{equation*}
\end{proof}

\subsection{Exercícios}

\begin{exercise}
Sejam $f$ e $g$ funções limitadas em $[a,b]$. Prove que
\begin{equation*}
\snorm{f+g}_{[a,b]} \leq \snorm{f}_{[a,b]} + \snorm{g}_{[a,b]} .
\end{equation*}
\end{exercise}

\begin{exercise}
\leavevmode
\begin{enumerate}[a)]
\item
Determine o limite pontual de
$\left\{ \dfrac{e^{x/n}}{n} \right\}_{n=1}^\infty$ para $x \in \R$.
\item
O limite é uniforme em $\R$?
\item
O limite é uniforme em $[0,1]$?
\end{enumerate}
\end{exercise}

\begin{exercise}
Suponha que as funções $f_n \colon S \to \R$ convirjam uniformemente
para $f \colon S \to \R$ e que $A \subset S$. Mostre que a sequência das
restrições $\{ f_n|_A \}_{n=1}^\infty$ converge uniformemente para $f|_A$.
\end{exercise}

\begin{exercise}
Suponha que $\{ f_n \}_{n=1}^\infty$ e $\{ g_n \}_{n=1}^\infty$, definidas
em um conjunto $A$, convirjam pontualmente para $f$ e $g$, respectivamente.
Mostre que $\{ f_n+g_n \}_{n=1}^\infty$ converge pontualmente para $f+g$.
\end{exercise}

\begin{exercise}
Suponha que $\{ f_n \}_{n=1}^\infty$ e $\{ g_n \}_{n=1}^\infty$, definidas
em um conjunto $A$, convirjam uniformemente em $A$ para $f$ e $g$,
respectivamente. Mostre que $\{ f_n+g_n \}_{n=1}^\infty$ converge
uniformemente em $A$ para $f+g$.
\end{exercise}

\begin{exercise}
Encontre um exemplo de um par de sequências de funções $\{ f_n \}_{n=1}^\infty$ e
$\{ g_n \}_{n=1}^\infty$
que convirjam uniformemente em algum conjunto $A$ para certas funções $f$
e $g$, mas para as quais $\{ f_ng_n \}_{n=1}^\infty$ (o produto) não
convirja uniformemente para $fg$ em $A$. Dica: tome $A \coloneqq \R$ e
$f(x)\coloneqq g(x) \coloneqq x$. Você pode até escolher $f_n = g_n$.
\end{exercise}

\begin{exercise}
Suponha que exista uma sequência de funções $\{ g_n \}_{n=1}^\infty$ que
convirja uniformemente para $0$ em $A$. Suponha também que tenhamos uma
sequência de funções $\{ f_n \}_{n=1}^\infty$ e uma função $f$ em $A$ tais que
\begin{equation*}
\babs{f_n(x) - f(x)} \leq g_n(x) 
\end{equation*}
para todo $x \in A$. Mostre que $\{ f_n \}_{n=1}^\infty$ converge
uniformemente para $f$ em $A$.
\end{exercise}

\begin{exercise}
Sejam $\{ f_n \}_{n=1}^\infty$, $\{ g_n \}_{n=1}^\infty$ e
$\{ h_n \}_{n=1}^\infty$ sequências de funções em $[a,b]$. Suponha que
$\{ f_n \}_{n=1}^\infty$ e $\{ h_n \}_{n=1}^\infty$ convirjam
uniformemente para uma função $f \colon [a,b] \to \R$ e que
$f_n(x) \leq g_n(x) \leq h_n(x)$ para todo $x \in [a,b]$. Mostre que
$\{ g_n \}_{n=1}^\infty$ converge uniformemente para $f$.
\end{exercise}

\begin{exercise}
Seja $f_n \colon [0,1] \to \R$ uma sequência de funções crescentes (isto
é, $f_n(x) \geq f_n(y)$ sempre que $x \geq y$). Suponha que $f_n(0) = 0$
e que $\lim\limits_{n \to \infty} f_n(1) = 0$. Mostre que
$\{ f_n \}_{n=1}^\infty$ converge uniformemente para $0$.
\end{exercise}

\begin{exercise}
Considere uma sequência de funções $f_n \colon [0,1] \to \R$ para a qual
existe uma sequência de números distintos $x_n \in [0,1]$ tal que, para
todo $n$,
\begin{equation*}
f_n(x_n) = 1.
\end{equation*}
Prove ou refute as afirmações seguintes:
\begin{enumerate}[a)]
\item
Verdadeiro ou falso: existe uma sequência $\{ f_n \}_{n=1}^\infty$ como
acima que converge pontualmente para $0$.
\item
Verdadeiro ou falso: existe uma sequência $\{ f_n \}_{n=1}^\infty$ como
acima que converge uniformemente para $0$.
\end{enumerate}
\end{exercise}

\begin{exercise}
Fixe uma função contínua $h \colon [a,b] \to \R$.
Defina $f(x) \coloneqq h(x)$ para $x \in [a,b]$,
$f(x) \coloneqq h(a)$ para $x < a$ e $f(x) \coloneqq h(b)$ para todo
$x > b$. Primeiro, mostre que $f \colon \R \to \R$ é contínua.
Agora, seja $f_n$ a função $g$ do
\exerciseref{exercise:smoothingout}, com $\epsilon = \nicefrac{1}{n}$,
definida no intervalo $[a,b]$. Isto é,
\begin{equation*}
f_n(x) \coloneqq \frac{n}{2} \int_{x-1/n}^{x+1/n} f .
\end{equation*}
Mostre que $\{ f_n \}_{n=1}^\infty$ converge uniformemente para $h$ em $[a,b]$.
\end{exercise}


\begin{exercise}
Prove que, se uma sequência de funções $f_n \colon S \to \R$ converge
uniformemente para uma função limitada $f \colon S \to \R$, então existe
um $N$ tal que $f_n$ é limitada para todo $n \geq N$.
\end{exercise}

\begin{exercise}
Suponha que exista uma constante $B$ tal que uma sequência de funções
$f_n \colon S \to \R$ seja limitada por $B$, isto é,
$\babs{f_n(x)} \leq B$ para todo $x \in S$.
Suponha que $\{ f_n \}_{n=1}^\infty$ convirja pontualmente para
$f \colon S \to \R$. Prove que $f$ é limitada.
\end{exercise}

\begin{exercise}[requer \sectionref{sec:moreonseries}]
No \exampleref{example:geomsumptconv}, vimos que
$\sum_{k=0}^\infty x^k$ converge pontualmente para $\frac{1}{1-x}$ em
$(-1,1)$.
\begin{enumerate}[a)]
\item
Mostre que, sempre que $0 \leq c < 1$, a série
$\sum_{k=0}^\infty x^k$ converge uniformemente em $[-c,c]$.
\item
Mostre que a série $\sum_{k=0}^\infty x^k$ não converge uniformemente
em $(-1,1)$.
\end{enumerate}
\end{exercise}

%%%%%%%%%%%%%%%%%%%%%%%%%%%%%%%%%%%%%%%%%%%%%%%%%%%%%%%%%%%%%%%%%%%%%%%%%%%%%%

\sectionnewpage
\section{Troca da ordem dos limites}
\label{sec:liminter}

%mbxINTROSUBSECTION

\sectionnotes{1--2,5 aulas; as subseções sobre derivadas e séries de
potências (que requerem a \sectionref{sec:moreonseries}) são opcionais.}

Boa parte da análise moderna trata principalmente da questão de trocar a
ordem de duas operações de limite. Quando temos uma cadeia de dois limites,
nem sempre podemos simplesmente trocar sua ordem. Por exemplo,
\begin{equation*}
0 = 
\lim_{n\to\infty}
\left(
\lim_{k\to\infty}
\frac{n}{n + k}
\right)
\neq
\lim_{k\to\infty}
\left(
\lim_{n\to\infty}
\frac{n}{n + k}
\right)
= 1 .
\end{equation*}

Ao tratar de sequências de funções, a troca da ordem dos limites aparece
com frequência. Vamos examinar vários casos: a continuidade, a integral e
a derivada do limite, e a convergência de séries de potências.

\subsection{Continuidade do limite}

Se temos uma sequência $\{ f_n \}_{n=1}^\infty$ de funções contínuas,
o limite também é contínuo? Suponha que $f$ seja o limite (pontual) de
$\{ f_n \}_{n=1}^\infty$. Se $\lim_{k\to\infty} x_k = x$, queremos
examinar a seguinte troca da ordem dos limites, em que a igualdade a provar
está marcada com um ponto de interrogação. Essa igualdade nem sempre vale.
Na verdade, talvez nem existam os limites à esquerda do ponto de
interrogação.
\begin{equation*}
\lim_{k \to \infty} 
f(x_k)
=
\lim_{k \to \infty} 
\Bigl(
\lim_{n \to \infty} f_n(x_k)
\Bigr)
\overset{?}{=}
\lim_{n \to \infty}
\Bigl(
\lim_{k \to \infty} 
f_n(x_k)
\Bigr)
=
\lim_{n \to \infty}
f_n(x)
=
f(x) .
\end{equation*}
Queremos encontrar condições para a sequência $\{ f_n \}_{n=1}^\infty$
que garantam a igualdade acima. Se exigirmos apenas convergência pontual,
o limite de uma sequência de funções não precisa ser contínuo, e a igualdade
acima pode não valer.

\begin{example}
Defina $f_n \colon [0,1] \to \R$ por
\begin{equation*}
f_n(x) \coloneqq
\begin{cases}
1-nx &  \text{se } x < \nicefrac{1}{n},\\
0 &  \text{se } x \geq \nicefrac{1}{n}.
\end{cases}
\end{equation*}
Veja a \figureref{contconvcntr:fig}.

\begin{myfigureht}
\myincludegraphics{contconvcntr}{%
Gráfico de uma função no intervalo de 0 a 1: ela vale 1 no extremo
esquerdo e, em seguida, seu gráfico é um segmento de reta que desce até 0
no eixo $x$, quando $x$ vale $1/n$. Depois, a função permanece nula até o
extremo direito.}
\caption{Gráfico de $f_n(x)$.%
\label{contconvcntr:fig}}
\end{myfigureht}

Cada função $f_n$ é contínua.
Fixe $x \in (0,1]$. Se $n \geq \nicefrac{1}{x}$, então
$x \geq \nicefrac{1}{n}$. Portanto, para $n \geq \nicefrac{1}{x}$,
temos $f_n(x) = 0$ e, assim,
\begin{equation*}
\lim_{n \to \infty} f_n(x) = 0.
\end{equation*}
Por outro lado, se $x=0$, então
\begin{equation*}
\lim_{n \to \infty} f_n(0) = 
\lim_{n \to \infty} 1 = 1.
\end{equation*}
Logo, o limite pontual de $f_n$ é a função
$f \colon [0,1] \to \R$ definida por
\begin{equation*}
f(x) \coloneqq
\begin{cases}
1 &  \text{se } x = 0,\\
0 &  \text{se } x > 0.
\end{cases}
\end{equation*}
A função $f$ não é contínua em 0.
\end{example}

Se, porém, exigirmos que a convergência seja uniforme, poderemos trocar a
ordem dos limites.

\begin{thm}
Suponha que $S \subset \R$ e que $\{ f_n \}_{n=1}^\infty$ seja uma
sequência de funções contínuas $f_n \colon S \to \R$ que converge
uniformemente para $f \colon S \to \R$. Então $f$ é contínua.
\end{thm}

\begin{proof}
Fixe $x \in S$. Seja $\{ x_n \}_{n=1}^\infty$ uma sequência em $S$ que
converge para $x$.

Seja $\epsilon > 0$ dado.
Como $\{ f_k \}_{k=1}^\infty$ converge uniformemente para $f$, existe
$k \in \N$ tal que
\begin{equation*}
\babs{f_k(y)-f(y)} < \nicefrac{\epsilon}{3}
\end{equation*}
para todo $y \in S$. Como $f_k$ é contínua em $x$, existe $N \in \N$ tal
que, para todo $m \geq N$,
\begin{equation*}
\babs{f_k(x_m)-f_k(x)} < \nicefrac{\epsilon}{3} .
\end{equation*}
Assim, para todo $m \geq N$,
\begin{equation*}
\begin{split}
\babs{f(x_m)-f(x)}
& =
\babs{f(x_m)-f_k(x_m)+f_k(x_m)-f_k(x)+f_k(x)-f(x)}
\\
& \leq
\babs{f(x_m)-f_k(x_m)}+
\babs{f_k(x_m)-f_k(x)}+
\babs{f_k(x)-f(x)}
\\
& <
\nicefrac{\epsilon}{3} +
\nicefrac{\epsilon}{3} +
\nicefrac{\epsilon}{3} = \epsilon .
\end{split}
\end{equation*}
Portanto, $\bigl\{ f(x_m) \bigr\}_{m=1}^\infty$ converge para $f(x)$ e,
consequentemente, $f$ é contínua em $x$.
Como $x$ foi escolhido arbitrariamente, $f$ é contínua em todo ponto.
\end{proof}

\subsection{Integral do limite}

Novamente, se exigirmos apenas convergência pontual, a integral do limite
de uma sequência de funções pode não ser igual ao limite das integrais.

\begin{example}
Defina $f_n \colon [0,1] \to \R$ por
\begin{equation*}
f_n(x) \coloneqq
\begin{cases}
0 &  \text{se } x = 0,\\
n-n^2x &  \text{se } 0 < x < \nicefrac{1}{n},\\
0 &  \text{se } x \geq \nicefrac{1}{n}.
\end{cases}
\avoidbreak
\end{equation*}
Veja a \figureref{intconvcntr:fig}.

\begin{myfigureht}
\myincludegraphics{intconvcntr}{%
Gráfico de uma função no intervalo de 0 a 1: ela vale 0 no extremo
esquerdo, mas então apresenta uma descontinuidade de salto e segue uma reta
descendente do ponto $(0,n)$ até o eixo $x$, onde $x=1/n$. Depois, vale 0
no restante do intervalo.}
\caption{Gráfico de $f_n(x)$.%
\label{intconvcntr:fig}}
\end{myfigureht}

Cada $f_n$ é integrável à Riemann (é contínua em $(0,1]$ e limitada), e o
teorema fundamental do cálculo nos dá
\begin{equation*}
\int_0^1 f_n =
\int_0^{\nicefrac{1}{n}} (n-n^2x)\,dx = \nicefrac{1}{2} .
\end{equation*}
Calculemos o limite pontual de $\{ f_n \}_{n=1}^\infty$.
Fixe $x \in (0,1]$. Para $n \geq \nicefrac{1}{x}$,
temos $x \geq \nicefrac{1}{n}$ e, portanto, $f_n(x) = 0$. Logo,
\begin{equation*}
\lim_{n \to \infty} f_n(x) = 0.
\end{equation*}
Além disso, $f_n(0) = 0$ para todo $n$. Assim, o limite pontual de
$\{ f_n \}_{n=1}^\infty$ é a função nula. Em resumo,
\begin{equation*}
\nicefrac{1}{2} =
\lim_{n\to\infty}
\int_0^1 f_n (x)\,dx
\neq
\int_0^1
\left(
\lim_{n\to\infty}
f_n(x)\right)\,dx
=
\int_0^1 0\,dx = 0 .
\end{equation*}
\end{example}

Mas, se exigirmos convergência uniforme, poderemos trocar a ordem dos
limites.\footnote{Condições mais fracas são suficientes para esse tipo de
teorema, mas provar tal generalização exige ferramentas mais sofisticadas
do que as que abordamos aqui: a integral de Lebesgue. Em particular, o
teorema vale para convergência pontual desde que $f$ seja integrável e
exista $M$ tal que $\snorm{f_n}_{[a,b]} \leq M$ para todo $n$.}

\begin{thm} \label{integralinterchange:thm}
Seja $\{ f_n \}_{n=1}^\infty$ uma sequência de funções
$f_n \colon [a,b] \to \R$ integráveis à Riemann que converge
uniformemente para $f \colon [a,b] \to \R$. Então $f$ é integrável à
Riemann e
\begin{equation*}
\int_a^b f = \lim_{n\to\infty} \int_a^b f_n .
\end{equation*}
\end{thm}

\begin{proof}
Seja $\epsilon > 0$ dado.
Como $f_n$ converge uniformemente para $f$, existe $M \in \N$ tal que,
para todo $n \geq M$, temos
$\babs{f_n(x)-f(x)} < \frac{\epsilon}{2(b-a)}$ para todo $x \in [a,b]$.
Em particular, pela desigualdade triangular reversa,
$\babs{f(x)} < \frac{\epsilon}{2(b-a)} + \babs{f_n(x)}$ para todo $x$.
Logo, $f$ é limitada, pois $f_n$ é limitada.
Como $f_n$ é integrável, calculemos
\begin{equation*}
\begin{split}
\overline{\int_a^b} f
-
\underline{\int_a^b} f
& =
\overline{\int_a^b} \bigl( f(x) - f_n(x) + f_n(x) \bigr)\,dx
-
\underline{\int_a^b} \bigl( f(x) - f_n(x) + f_n(x) \bigr)\,dx
\\
& \leq
\overline{\int_a^b} \bigl( f(x) - f_n(x) \bigr)\,dx +  \overline{\int_a^b} f_n(x) \,dx
-
\underline{\int_a^b} \bigl( f(x) - f_n(x) \bigr)\,dx -  \underline{\int_a^b} f_n(x) \,dx
\\
& =
\overline{\int_a^b} \bigl( f(x) - f_n(x) \bigr)\,dx +  \int_a^b f_n(x) \,dx
-
\underline{\int_a^b} \bigl( f(x) - f_n(x) \bigr)\,dx -  \int_a^b f_n(x) \,dx
\\
& =
\overline{\int_a^b} \bigl( f(x) - f_n(x) \bigr)\,dx
-
\underline{\int_a^b} \bigl( f(x) - f_n(x) \bigr)\,dx
\\
& \leq
\frac{\epsilon}{2(b-a)} (b-a) + 
\frac{\epsilon}{2(b-a)} (b-a) = \epsilon .
\end{split}
\end{equation*}
Na primeira desigualdade, usamos a \propref{prop:upperlowerlinineq}.
A segunda decorre da \propref{intulbound:prop} e do fato de que, para
todo $x \in [a,b]$,
$\frac{-\epsilon}{2(b-a)} < f(x)-f_n(x) < \frac{\epsilon}{2(b-a)}$.
Como $\epsilon > 0$ foi arbitrário, $f$ é integrável à Riemann.

Por fim, calculemos $\int_a^b f$. Aplicamos a \propref{intbound:prop}
na conta. Novamente, para todo $n \geq M$ (com o mesmo $M$ de antes),
\begin{equation*}
\begin{split}
\abs{\int_a^b f - \int_a^b f_n} & = 
\abs{ \int_a^b \bigl(f(x) - f_n(x)\bigr)\,dx}
\\
& \leq
\frac{\epsilon}{2(b-a)} (b-a) = \frac{\epsilon}{2} < \epsilon .
\end{split}
\end{equation*}
Portanto, $\bigl\{ \int_a^b f_n \bigr\}_{n=1}^\infty$ converge para
$\int_a^b f$.
\end{proof}

\begin{example}
Suponha que queiramos calcular
\begin{equation*}
\lim_{n\to\infty} \int_0^1 \frac{nx+ \sin(nx^2)}{n} \,dx .
\end{equation*}
É impossível calcular as integrais para um $n$ específico usando o
cálculo, pois $\sin(nx^2)$ não tem uma primitiva expressável por funções
elementares. No entanto, podemos calcular o limite.
Já mostramos que $\frac{nx+ \sin(nx^2)}{n}$ converge uniformemente em
$[0,1]$ para $x$.
Pelo \thmref{integralinterchange:thm}, o limite existe e
\begin{equation*}
\lim_{n\to\infty} \int_0^1 \frac{nx+ \sin(nx^2)}{n} \,dx
=
\int_0^1
x \,dx = \nicefrac{1}{2} .
\end{equation*}
\end{example}

\begin{example}
Se a convergência for apenas pontual, o limite pode nem sequer ser
integrável à Riemann. Em $[0,1]$, defina
\begin{equation*}
f_n(x) \coloneqq
\begin{cases}
1 & \text{se } x = \nicefrac{p}{q} \text{ na forma irredutível e } q \leq n, \\
0 & \text{caso contrário.}
\end{cases}
\end{equation*}
Cada função $f_n$ difere da função nula em apenas um número finito de
pontos, pois há somente um número finito de frações em $[0,1]$ cujo
denominador é menor ou igual a $n$. Portanto, $f_n$ é integrável e
$\int_0^1 f_n = \int_0^1 0 = 0$. É fácil mostrar que
$\{ f_n \}_{n=1}^\infty$ converge pontualmente para a
\myindex{função de Dirichlet}
\begin{equation*}
f(x) \coloneqq
\begin{cases}
1 & \text{se } x \in \Q, \\
0 & \text{caso contrário,}
\end{cases}
\end{equation*}
que não é integrável à Riemann.
\end{example}

\begin{example}
Na verdade, se a convergência for apenas pontual, o limite de funções
limitadas pode nem sequer ser limitado.
Defina $f_n \colon [0,1] \to \R$ por
\begin{equation*}
f_n(x) \coloneqq
\begin{cases}
0 & \text{se } x < \nicefrac{1}{n},\\
\nicefrac{1}{x} & \text{caso contrário.}
\end{cases}
\end{equation*}
Para todo $n$, temos $\babs{f_n(x)} \leq n$ para todo $x \in [0,1]$,
portanto as funções são limitadas. No entanto,
$\{ f_n \}_{n=1}^\infty$ converge pontualmente para a função ilimitada
\begin{equation*}
f(x) \coloneqq
\begin{cases}
0 & \text{se } x = 0,\\
\nicefrac{1}{x} & \text{caso contrário.}
\end{cases}
\end{equation*}
\end{example}

\subsection{Derivada do limite}

Embora a convergência uniforme seja suficiente para trocar a ordem dos
limites com integrais, ela não basta para fazer o mesmo com derivadas, a
menos que as próprias derivadas também convirjam uniformemente.

\begin{example}
Seja $f_n(x) \coloneqq \frac{\sin(nx)}{n}$. Então $f_n$ converge
uniformemente para 0. Veja a \figureref{fig:conv1nsinxn}.
A derivada do limite é 0. No entanto, $f_n'(x) = \cos(nx)$ não converge
nem sequer pontualmente.
Por exemplo, $f_n'(\pi) = {(-1)}^n$.
Além disso, $f_n'(0) = 1$ para todo $n$; essa sequência converge, mas não
para 0.
\begin{myfigureht}
\myincludegraphics{conv1nsinxn}{%
Ondas senoidais com 10 frequências diferentes, agora também reescaladas
para baixo nas frequências mais altas. As funções de frequência maior ficam
mais próximas do eixo $x$; por isso, os gráficos já não preenchem uma faixa
inteira.}
\caption{Gráficos de $\frac{\sin(nx)}{n}$ para
$n=1,2,\ldots,10$, com os valores maiores de $n$ em tons mais claros de cinza.%
\label{fig:conv1nsinxn}}
\end{myfigureht}
\end{example}

\begin{example} \label{exercise:badconvergenceder}
Seja $f_n(x) \coloneqq \frac{1}{1+nx^2}$.
Se $x \neq 0$, então $\lim_{n \to \infty} f_n(x) = 0$,
mas $\lim_{n \to \infty} f_n(0) = 1$.
Logo, $\{ f_n \}_{n=1}^\infty$ converge pontualmente para uma função que
não é contínua em 0.
Calculamos
\begin{equation*}
f_n'(x) %= \frac{d}{dx} \frac{1}{1+n x^2}
= \frac{-2 n x}{(1+ n x^2)^2} .
\end{equation*}
Para todo $x$, $\lim_{n\to\infty} f_n'(x) = 0$, portanto as derivadas
convergem pontualmente para 0, mas o leitor pode verificar que a
convergência não é uniforme em nenhum intervalo que contenha 0.
O limite de $f_n$ não é diferenciável em 0 — nem sequer é contínuo em 0.
Veja a \figureref{fig:convder1over1pnxsq}.
\begin{myfigureht}
\myincludegraphics{convder1over1pnxsq}{%
Dois gráficos, cada um com várias funções em diferentes tons de cinza.
À esquerda, uma curva em forma de sino, cujo pico fica mais estreito, mas
mantém a mesma altura, nos tons mais claros. À direita, a derivada: ela é
positiva e tem um pico antes da origem; depois passa pela origem, atinge um
mínimo após a origem e volta a se aproximar do eixo $x$. À medida que o tom
fica mais claro, os picos e mínimos ficam mais estreitos e se afastam do
eixo $x$.}
\caption{Gráficos de $\frac{1}{1+nx^2}$ e de sua derivada para
$n=1,2,\ldots,10$, com os valores maiores de $n$ em tons mais claros de cinza.%
\label{fig:convder1over1pnxsq}}
\end{myfigureht}
\end{example}

Veja mais exemplos nos exercícios. Usando o teorema fundamental do cálculo,
obtemos um resultado para funções continuamente diferenciáveis. O teorema
seguinte também vale sem supor que as derivadas sejam contínuas, mas sua
demonstração é mais difícil.

\begin{thm} \label{thm:dersconverge}
Seja $I$ um intervalo limitado e sejam $f_n \colon I \to \R$ funções
continuamente diferenciáveis.
Suponha que $\{ f_n' \}_{n=1}^\infty$ convirja uniformemente para
$g \colon I \to \R$ e que $\bigl\{ f_n(c) \bigr\}_{n=1}^\infty$ seja
uma sequência convergente para algum $c \in I$. Então
$\{ f_n \}_{n=1}^\infty$ converge uniformemente para uma função
continuamente diferenciável $f \colon I \to \R$, e $f' = g$.
\end{thm}

\begin{proof}
Defina $f(c) \coloneqq \lim_{n\to \infty} f_n(c)$.
Como $f_n'$ é contínua e, portanto, integrável à Riemann, o teorema
fundamental do cálculo nos dá, para $x \in I$,
\begin{equation*}
f_n(x) = f_n(c) + \int_c^x f_n' .
\end{equation*}
Como $\{ f_n' \}_{n=1}^\infty$ converge uniformemente em $I$, também
converge uniformemente em $[c,x]$ (ou em $[x,c]$, se $x < c$).
Assim, o limite do lado direito quando $n \to \infty$ existe.
Defina $f$ nos demais pontos (isto é, quando $x\neq c$) por esse limite:
\begin{equation*}
f(x) \coloneqq
\lim_{n\to\infty} f_n(c) + \lim_{n\to\infty} \int_c^x f_n'
=
f(c) + \int_c^x g .
\end{equation*}
Como limite uniforme de funções contínuas, $g$ é contínua. Logo, pela
segunda forma do teorema fundamental, $f$ é diferenciável e
$f'(x) = g(x)$ para todo $x \in I$.

Resta provar a convergência uniforme.
Suponha que $I$ tenha extremo inferior $a$ e extremo superior $b$.
Seja $\epsilon > 0$ dado. Escolha $M$ tal que, para todo $n \geq M$,
$\babs{f(c)-f_n(c)} < \nicefrac{\epsilon}{2}$
e
$\babs{g(x)-f_n'(x)} < \frac{\epsilon}{2(b-a)}$
para todo $x \in I$. Então
\begin{equation*}
\begin{split}
\babs{f(x) - f_n(x)} & =
\abs{\left(f(c) + \int_c^x g\right) - \left( f_n(c) + \int_c^x f_n' \right)}
\\
& \leq
\babs{f(c) - f_n(c)} + \abs{\int_c^x g - \int_c^x f_n'}
\\
& =
\babs{f(c) - f_n(c)} + \abs{\int_c^x \bigl(g(s) - f_n'(s)\bigr) \, ds}
\\
& <
\frac{\epsilon}{2}
+
\frac{\epsilon}{2(b-a)}
(b-a)
=\epsilon. \qedhere
\end{split}
\end{equation*}
\end{proof}

O restante da demonstração vale mesmo quando $I$ não é limitado; a
limitação só é necessária para a convergência uniforme de $f_n$ para $f$.
Por exemplo, suponha que $I = \R$ e tome $f_n(x) \coloneqq \nicefrac{x}{n}$.
Então $f_n'(x)=\nicefrac{1}{n}$ converge uniformemente para 0. No entanto,
$\{f_n\}_{n=1}^\infty$ converge para 0 apenas pontualmente.

\subsection{Convergência de séries de potências}

Na \sectionref{sec:moreonseries}, vimos que uma série de potências converge
absolutamente no interior de seu raio de convergência e, portanto, converge
pontualmente. Vamos mostrar que ela — e também todas as suas derivadas —
converge uniformemente. Isso nos permite trocar a ordem de vários tipos de
limite. Além de ser contínua, a função limite também pode ser integrada e
até diferenciada termo a termo.

\begin{prop}
Seja $\sum_{n=0}^\infty c_n {(x-a)}^n$ uma série de potências convergente
com raio de convergência $\rho$, em que $0 < \rho \leq \infty$.
Então a série converge uniformemente em $[a-r,a+r]$ sempre que
$0 < r < \rho$.

Em particular, a série converge pontualmente para uma função contínua em
$(a-\rho,a+\rho)$ se $\rho < \infty$, ou em $\R$ se $\rho = \infty$.
\end{prop}

\begin{proof}
Defina $I \coloneqq (a-\rho,a+\rho)$ se $\rho < \infty$,
ou $I \coloneqq \R$ se $\rho= \infty$.
Tome $0 < r < \rho$.
A série converge absolutamente para todo $x \in I$, em particular para
$x = a+r$. Logo, $\sum_{n=0}^\infty \sabs{c_n} r^n$ converge.
Dado $\epsilon >0$, escolha $M$ tal que, para todo $k \geq M$,
\begin{equation*}
\sum_{n=k+1}^\infty \sabs{c_n} {r}^n < \epsilon .
\end{equation*}
Para todo $x \in [a-r,a+r]$ e todo $m > k$,
\begin{multline*}
\abs{\sum_{n=0}^m c_n {(x-a)}^n - 
\sum_{n=0}^k c_n {(x-a)}^n}
=
\abs{\sum_{n=k+1}^m c_n {(x-a)}^n}
\\
\leq
\sum_{n=k+1}^m \sabs{c_n} {\sabs{x-a}}^n
\leq
\sum_{n=k+1}^m \sabs{c_n} {r}^n
\leq
\sum_{n=k+1}^\infty \sabs{c_n} {r}^n
<\epsilon.
\end{multline*}
Portanto, as somas parciais formam uma sequência uniformemente de Cauchy
em $[a-r,a+r]$ e, assim, convergem uniformemente nesse conjunto.

Além disso, as somas parciais são polinômios e, portanto, contínuas. Assim,
seu limite uniforme em $[a-r,a+r]$ é uma função contínua.
Como $r < \rho$ foi arbitrário, a função limite é contínua em todo $I$.
\end{proof}

Como dissemos, mostraremos que séries de potências podem ser diferenciadas
e integradas termo a termo. A série obtida por diferenciação ou integração
também é uma série de potências, e mostraremos que ela tem o mesmo raio de
convergência. Portanto, toda série de potências define uma função
infinitamente diferenciável.

Primeiro provaremos que podemos integrar termo a termo, pois a integração
exige apenas convergência uniforme.

\begin{cor}
Seja $\sum_{n=0}^\infty c_n {(x-a)}^n$ uma série de potências convergente
com raio de convergência $0 < \rho \leq \infty$.
Defina $I \coloneqq (a-\rho,a+\rho)$ se $\rho < \infty$,
ou $I \coloneqq \R$ se $\rho= \infty$, e seja $f \colon I \to \R$ o
limite da série. Então
\begin{equation*}
\int_a^x f = \sum_{n=1}^\infty \frac{c_{n-1}}{n} {(x-a)}^{n} ,
\end{equation*}
e o raio de convergência dessa série é pelo menos $\rho$.
\end{cor}

\begin{proof}
Tome $0 < r < \rho$.
As somas parciais $\sum_{n=0}^k c_n {(x-a)}^n$ convergem uniformemente em
$[a-r,a+r]$.
Para cada $x \in [a-r,a+r]$ fixado, a convergência também é uniforme em
$[a,x]$ (ou em $[x,a]$, se $x < a$).
Logo,
\begin{equation*}
\int_a^x f =
\int_a^x \lim_{k\to\infty} \sum_{n=0}^k c_n {(s-a)}^n \, ds
=
\lim_{k\to\infty}
\int_a^x \sum_{n=0}^k c_n {(s-a)}^n \, ds
=
\lim_{k\to\infty}
\sum_{n=1}^{k+1} \frac{c_{n-1}}{n} {(x-a)}^{n} . \qedhere
\end{equation*}
\end{proof}


\begin{cor} \label{cor:differentiatepowerser}
Seja $\sum_{n=0}^\infty c_n {(x-a)}^n$ uma série de potências convergente
com raio de convergência $0 < \rho \leq \infty$.
Defina $I \coloneqq (a-\rho,a+\rho)$ se $\rho < \infty$,
ou $I \coloneqq \R$ se $\rho= \infty$, e seja $f \colon I \to \R$ o
limite da série. Então $f$ é diferenciável e
\begin{equation*}
f'(x) = \sum_{n=0}^\infty (n+1) c_{n+1} {(x-a)}^{n} ,
\end{equation*}
e o raio de convergência dessa série é $\rho$.
\end{cor}

\begin{proof}
Tome $0 < r < \rho$.
A série converge uniformemente em $[a-r,a+r]$, mas precisamos da
convergência uniforme da derivada.
Defina
\begin{equation*}
R \coloneqq \limsup_{n \to \infty} \sabs{c_n}^{1/n} .
\end{equation*}
Como a série converge, $R < \infty$.
O raio de convergência $\rho$ é $\nicefrac{1}{R}$ (ou $\rho = \infty$ se
$R=0$).

Seja $\epsilon > 0$ dado. No \exampleref{example:nto1overn}, vimos que
$\lim_{n\to\infty} n^{1/n} = 1$.
Portanto, existe $N$ tal que, para todo $n \geq N$,
$n^{1/n} < 1+\epsilon$.
Assim,
\begin{equation*}
R = 
\limsup_{n \to \infty}
\sabs{c_n}^{1/n}
\leq
\limsup_{n \to \infty}
\sabs{n c_n}^{1/n}
\leq
(1+\epsilon)
\limsup_{n \to \infty}
\sabs{c_n}^{1/n}
=
(1+\epsilon)R .
\end{equation*}
Como $\epsilon$ foi arbitrário,
$\limsup_{n \to \infty} \sabs{n c_n}^{1/n} = R$.
Portanto, $\sum_{n=1}^\infty n c_{n} {(x-a)}^{n}$ tem raio de
convergência $\rho$. Dividindo por $(x-a)$, concluímos que
$\sum_{n=0}^\infty (n+1) c_{n+1} {(x-a)}^{n}$ também tem raio de
convergência $\rho$.

Consequentemente, as somas parciais
$\sum_{n=0}^k (n+1) c_{n+1} {(x-a)}^{n}$,
que são as derivadas das somas parciais
$\sum_{n=0}^{k+1} c_{n} {(x-a)}^{n}$,
convergem uniformemente em $[a-r,a+r]$. Além disso, a série claramente
converge em $x=a$.
Podemos, então, aplicar o \thmref{thm:dersconverge}. Isso conclui a prova,
pois $r < \rho$ foi arbitrário.
\end{proof}

\begin{example} \label{example:exponentialbypowerseries}
Poderíamos ter usado esse resultado para definir a função exponencial. Isto
é, a série de potências
\begin{equation*}
f(x) \coloneqq \sum_{n=0}^\infty \frac{x^n}{n!}
\end{equation*}
tem raio de convergência $\rho=\infty$. Além disso, $f(0) = 1$.
Derivando termo a termo, obtemos $f'(x) = f(x)$.
\end{example}

\begin{example}
A série
\begin{equation*}
\sum_{n=1}^\infty n x^n
\end{equation*}
converge para $\frac{x}{{(1-x)}^2}$ em $(-1,1)$.

Demonstração:
Em $(-1,1)$, $\sum_{n=0}^\infty x^n$ converge para
$\frac{1}{1-x}$. A série derivada $\sum_{n=1}^\infty n x^{n-1}$ converge,
então, no mesmo intervalo para $\frac{1}{{(1-x)}^2}$. Multiplicando o
resultado por $x$, obtemos o que queríamos.
\end{example}

\subsection{Exercícios}

\begin{exercise}
%Embora a convergência uniforme preserve a continuidade, ela não preserva
%a diferenciabilidade.
Encontre um exemplo explícito de uma sequência de funções diferenciáveis
em $[-1,1]$ que convirja uniformemente para uma função $f$ que não seja
diferenciável.
Dica: há muitas possibilidades. Uma das mais simples talvez seja combinar
$\sabs{x}$ com $\frac{n}{2}x^2 + \frac{1}{2n}$.
Outra possibilidade é considerar $\sqrt{x^2+{(\nicefrac{1}{n})}^2}$.
Mostre que essas funções são diferenciáveis, que convergem uniformemente e,
em seguida, que o limite não é diferenciável.
\end{exercise}

\begin{exercise}
Seja $f_n(x) \coloneqq \frac{x^n}{n}$. Mostre que
$\{ f_n \}_{n=1}^\infty$ converge uniformemente para uma função
diferenciável $f$ em $[0,1]$ (determine $f$). Mostre, porém, que
$f'(1) \neq \lim\limits_{n\to\infty} f_n'(1)$.
\end{exercise}

%\begin{exnote}
%Nota: os dois exercícios anteriores mostram que não podemos simplesmente
%trocar a ordem dos limites com derivadas, mesmo quando a convergência é
%uniforme. Veja também o exercício \exerciseref{c1uniflim:exercise} abaixo.
%\end{exnote}
%
\begin{exercise}
Seja $f \colon [0,1] \to \R$ integrável à Riemann (portanto, limitada).
Determine $\displaystyle \lim_{n\to\infty} \int_0^1 \frac{f(x)}{n} \,dx$.
\end{exercise}

\begin{exercise}
Mostre que
$\displaystyle \lim_{n\to\infty} \int_1^2 e^{-nx^2} \,dx = 0$. Você pode
usar fatos sobre a função exponencial conhecidos do cálculo.
%o que você sabe sobre a função exponencial do cálculo.
\end{exercise}

\begin{exercise}
Encontre um exemplo de uma sequência de funções contínuas em $(0,1)$ que
converge pontualmente para uma função contínua em $(0,1)$, mas não converge
uniformemente.
\end{exercise}

\begin{exnote}
Nota: no exercício anterior, escolhemos $(0,1)$ por simplicidade. Para um
exercício mais desafiador, substitua $(0,1)$ por $[0,1]$.
\end{exnote}

\begin{exercise}
Decida se a afirmação seguinte é verdadeira ou falsa; prove-a ou encontre
um contraexemplo: se $\{ f_n \}_{n=1}^\infty$ é uma sequência de funções
descontínuas em todos os pontos de $[0,1]$ que converge uniformemente para
uma função $f$, então $f$ é descontínua em todos os pontos.
\end{exercise}

\begin{exercise} \label{c1uniflim:exercise}
Para uma função continuamente diferenciável $f \colon [a,b] \to \R$,
defina
\begin{equation*}
\snorm{f}_{C^1} \coloneqq \snorm{f}_{[a,b]} + \snorm{f'}_{[a,b]} .
\end{equation*}
Suponha que $\{ f_n \}_{n=1}^\infty$ seja uma sequência de funções
continuamente diferenciáveis tal que, para todo $\epsilon >0$, exista $M$
com a propriedade de que, para todos $n,k \geq M$,
\begin{equation*}
\snorm{f_n-f_k}_{C^1} < \epsilon .
\end{equation*}
Mostre que $\{ f_n \}_{n=1}^\infty$ converge uniformemente para alguma
função continuamente diferenciável $f \colon [a,b] \to \R$.
\end{exercise}

\begin{exnote}
Suponha que $f \colon [0,1] \to \R$ seja integrável à Riemann.
Para os dois exercícios seguintes, defina o número
\begin{equation*}
\snorm{f}_{L^1} \coloneqq 
\int_0^1 \babs{f(x)}\,dx .
\end{equation*}
De fato, $\sabs{f}$ é integrável sempre que $f$ é; veja o exercício
\exerciseref{exercise:hardabsint}.
Esse número é chamado de \emph{norma $L^1$}\index{L1-norma@$L^1$-norma} e
define outro tipo muito comum de convergência, chamado
\emph{convergência em $L^1$}\index{convergência em L1@$L^1$-convergência}.
Esse conceito, porém, é um pouco mais sutil.
\end{exnote}

\begin{exercise}
Suponha que $\{ f_n \}_{n=1}^\infty$ seja uma sequência de funções
integráveis à Riemann em $[0,1]$ que converge uniformemente para $0$.
Mostre que
\begin{equation*}
\lim_{n\to\infty} \snorm{f_n}_{L^1} = 0 .
\end{equation*}
\end{exercise}

\begin{exercise}
Encontre uma sequência $\{ f_n \}_{n=1}^\infty$ de funções integráveis à
Riemann em $[0,1]$ que convirja pontualmente para $0$, mas para a qual
\begin{equation*}
\lim_{n\to\infty} \snorm{f_n}_{L^1} \text{não existe (é } \infty\text{)}.
\end{equation*}
\end{exercise}

\begin{exercise}[Difícil] \label{exercise:dinisthm}
Prove o \emph{\myindex{teorema de Dini}}:
Seja $f_n \colon [a,b] \to \R$ uma sequência de funções contínuas tal que
\begin{equation*}
0 \leq f_{n+1}(x) \leq f_n(x) \leq \cdots \leq f_1(x) 
\qquad \text{para todo } n \in \N.
\end{equation*}
Suponha que $\{ f_n \}_{n=1}^\infty$ convirja pontualmente para $0$.
Mostre que $\{ f_n \}_{n=1}^\infty$ converge uniformemente para zero.
\end{exercise}

\begin{exercise}
Suponha que $f_n \colon [a,b] \to \R$ seja uma sequência de funções
contínuas que converge pontualmente para uma função contínua
$f \colon [a,b] \to \R$. Suponha que, para todo $x \in [a,b]$, a
sequência $\bigl\{ \babs{f_n(x)-f(x)} \bigr\}_{n=1}^\infty$ seja
monótona. Mostre que $\{f_n\}_{n=1}^\infty$ converge uniformemente.
\end{exercise}

\begin{exercise}
\pagebreak[3]
Encontre sequências de funções integráveis à Riemann
$f_n \colon [0,1] \to \R$ tais que $\{ f_n \}_{n=1}^\infty$ convirja
pontualmente para zero e
\begin{enumerate}[a)]
\item
$\bigl\{ \int_0^1 f_n \bigr\}_{n=1}^\infty$ cresça sem limite,
\item
$\bigl\{ \int_0^1 f_n \bigr\}_{n=1}^\infty$ seja a sequência $-1,1,-1,1,-1,1, \ldots$.
\end{enumerate}
\end{exercise}

\begin{exnote}
É possível definir um
\emph{\myindex{limite conjunto}} de uma sequência dupla
$\{ x_{n,m} \}_{n,m=1}^\infty$ de números reais (isto é, de uma função
de $\N \times \N$ em $\R$).
Dizemos que $L$ é o limite conjunto de $\{ x_{n,m} \}_{n,m=1}^\infty$ e escrevemos
\begin{equation*}
\lim_{\substack{n\to\infty\\m\to\infty}}
x_{n,m} = L ,
\qquad
\text{ou}
\qquad
\lim_{(n,m) \to \infty}
x_{n,m} = L ,
\end{equation*}
se, para todo $\epsilon > 0$, existe $M$ tal que, se $n \geq M$ e
$m \geq M$, então
$\sabs{x_{n,m} - L} < \epsilon$.
\end{exnote}

\begin{exercise}
Suponha que o limite conjunto (definido acima) de
$\{ x_{n,m} \}_{n,m=1}^\infty$ seja $L$, e que, para todo $n$,
$\lim\limits_{m \to \infty} x_{n,m}$ exista e, para todo $m$,
$\lim\limits_{n \to \infty} x_{n,m}$ exista. Mostre então que
$\lim\limits_{n\to\infty}\lim\limits_{m \to \infty} x_{n,m}
=
\lim\limits_{m\to\infty}\lim\limits_{n \to \infty} x_{n,m} = L$.
\end{exercise}

\begin{exercise}
Ter um limite conjunto (definido acima) não significa que os limites iterados existam.
Considere $x_{n,m} \coloneqq \frac{{(-1)}^{n+m}}{\min \{n,m \}}$.
\begin{enumerate}[a)]
\item
Mostre que os limites $\lim\limits_{m \to \infty} x_{n,m}$ não existem
para nenhum $n$ e que os limites $\lim\limits_{n \to \infty} x_{n,m}$
não existem para nenhum $m$. Portanto, nenhum dos limites iterados
$\lim\limits_{n\to\infty}\lim\limits_{m \to \infty} x_{n,m}$ e
$\lim\limits_{m\to\infty}\lim\limits_{n \to \infty} x_{n,m}$ está definido.
\item
Mostre que o limite conjunto de $\{ x_{n,m} \}_{n,m=1}^\infty$ existe e é igual a 0.
\end{enumerate}
\end{exercise}

\begin{exercise}
Dizemos que uma sequência de funções $f_n \colon \R \to \R$
\emph{\myindex{converge uniformemente nos subconjuntos compactos}}
\index{convergência uniforme nos subconjuntos compactos}
se, para todo $k \in \N$, a sequência $\{ f_n \}_{n=1}^\infty$
converge uniformemente em $[-k,k]$.
\begin{enumerate}[a)]
\item
Prove que, se $f_n \colon \R \to \R$ é uma sequência de funções
contínuas que converge uniformemente nos subconjuntos compactos, então o
limite é contínuo.
\item 
Prove que, se $f_n \colon \R \to \R$ é uma sequência de funções
integráveis à Riemann em todo intervalo fechado e limitado $[a,b]$ e que
converge uniformemente nos subconjuntos compactos para uma função
$f \colon \R \to \R$, então, para todo intervalo $[a,b]$,
$f \in \sR\bigl([a,b]\bigr)$ e
$\int_a^b f = \lim\limits_{n\to\infty} \int_a^b f_n$.
\end{enumerate}
\end{exercise}

\begin{exercise}[Desafiador]
Encontre uma sequência de funções contínuas $f_n \colon [0,1] \to \R$
que convirja para a \myindex{função pipoca} $f \colon [0,1] \to \R$,
definida por $f(\nicefrac{p}{q}) \coloneqq \nicefrac{1}{q}$ (quando
$\nicefrac{p}{q}$ está na forma irredutível) e $f(x) \coloneqq 0$ quando
$x$ é irracional (note que $f(0) = f(1) = 1$); veja o
\exampleref{popcornfunction:example}.
Assim, o limite pontual de funções contínuas pode ter um conjunto denso de
descontinuidades. Veja também o exercício seguinte.
\end{exercise}

\begin{exercise}[Desafiador]
A \myindex{função de Dirichlet}
$f \colon [0,1] \to \R$, definida por $f(x) \coloneqq 1$ se $x \in \Q$
e $f(x) \coloneqq 0$ se $x \notin \Q$, não é limite pontual de funções
contínuas, embora isso seja difícil de provar.
Prove, porém, que $f$ é limite pontual de funções que, por sua vez, são
limites pontuais de funções contínuas.
\end{exercise}

\begin{exercise}
\leavevmode
\begin{enumerate}[a)]
\item
Encontre uma sequência de funções Lipschitz contínuas em $[0,1]$ cujo
limite uniforme seja $\sqrt{x}$, que não é uma função Lipschitz.
\item
Por outro lado, mostre que, se $f_n \colon S \to \R$ são funções Lipschitz
com uma constante uniforme $K$ (isto é, todas satisfazem a definição com
a mesma constante) e $\{ f_n \}_{n=1}^\infty$ converge pontualmente para
$f \colon S \to \R$, então o limite $f$ é uma função Lipschitz contínua
com constante de Lipschitz $K$.
\end{enumerate}
\end{exercise}

\begin{exercise}[requer \sectionref{sec:moreonseries}]
Se $\sum_{n=0}^\infty c_n {(x-a)}^n$ tem raio de convergência $\rho$,
mostre que a integral termo a termo
$\sum_{n=1}^\infty \frac{c_{n-1}}{n} {(x-a)}^n$ tem raio de convergência
$\rho$. Observe que, acima, provamos apenas que o raio de convergência é
pelo menos $\rho$.
\end{exercise}

\begin{exercise}[requer \sectionref{sec:moreonseries} e \sectionref{sec:taylor}]
\pagebreak[3]
Suponha que $f(x) \coloneqq \sum_{n=0}^\infty c_n {(x-a)}^n$ convirja em $(a-\rho,a+\rho)$.
\begin{enumerate}[a)]
\item
Suponha que $f^{(k)}(a) = 0$ para todo $k=0,1,2,3,\ldots$. Prove que
$c_n = 0$ para todo $n$ ou, em outras palavras, que
$f(x) = 0$ para todo $x \in (a-\rho,a+\rho)$.
\item
Usando o item a), prove uma versão do chamado
\myquote{teorema da identidade para funções analíticas}: se existe
$\epsilon > 0$ tal que $f(x) = 0$ para todo $x \in (a-\epsilon, a+\epsilon)$,
então $f(x) = 0$ para todo $x \in (a-\rho,a+\rho)$.
\end{enumerate}
\end{exercise}

\begin{exercise}
Seja $f_n(x) \coloneqq \frac{x}{1+{(nx)}^2}$. Observe que $f_n$ são
funções diferenciáveis.
\begin{enumerate}[a)]
\item
Mostre que $\{ f_n \}_{n=1}^\infty$ converge uniformemente para 0.
\item
Mostre que $\sabs{f_n'(x)} \leq 1$ para todo $x$ e todo $n$.
\item
Mostre que $\{ f_n' \}_{n=1}^\infty$ converge pontualmente para uma
função descontínua na origem.
\item
Seja $\{ a_n \}_{n=1}^\infty$ uma enumeração dos números racionais.
Defina
\begin{equation*}
g_n(x) \coloneqq \sum_{k=1}^n 2^{-k} f_n(x-a_k) .
\end{equation*}
Mostre que $\{ g_n \}_{n=1}^\infty$ converge uniformemente para 0.
\item
Mostre que $\{ g_n' \}_{n=1}^\infty$ converge pontualmente para uma função
$\psi$ que é descontínua em todo número racional e contínua em todo
irracional. Em particular, $\lim\limits_{n\to\infty} g_n'(x) \neq 0$ para
todo número racional $x$.
\end{enumerate}
\end{exercise}

\begin{exercise}[requer \sectionref{sec:impropriemann}]
Mostre que a convergência uniforme não basta para passar o limite através
de integrais impróprias em intervalos infinitos.
Isto é, encontre uma sequência de funções $f_n \colon \R \to \R$,
integráveis à Riemann em todo intervalo limitado, que convirja
uniformemente para zero e satisfaça $\int_{-\infty}^\infty f_n = 1$ para
todo $n$.
\end{exercise}

%%%%%%%%%%%%%%%%%%%%%%%%%%%%%%%%%%%%%%%%%%%%%%%%%%%%%%%%%%%%%%%%%%%%%%%%%%%%%%

\sectionnewpage
\section{Picard's theorem}
\label{sec:picard}

%mbxINTROSUBSECTION

\sectionnotes{1--2 lectures (can be safely skipped)}

A first semester course in analysis should have
a \emph{pi\`ece de r\'esistance} caliber
theorem.  We pick a theorem whose proof combines everything we have
learned.  It is more sophisticated than the fundamental theorem of calculus,
the first highlight theorem of this course.  The
theorem we are talking about is Picard's
theorem%
\footnote{Named for the French mathematician
\href{https://en.wikipedia.org/wiki/\%C3\%89mile_Picard}{Charles \'Emile Picard}
(1856--1941).}
on existence and uniqueness of a solution to an ordinary differential equation.
Both the statement and the proof are beautiful examples of what one can do
with the material we have mastered so far.  It is also a good example of how analysis is
applied, as differential equations are indispensable in science of every
stripe.

\subsection{First order ordinary differential equation}

Modern science is described in the language of
\emph{differential equations}\index{differential equation}.
That is, equations involving not only the unknown, but also its
derivatives.  The simplest nontrivial form of a differential equation is
the \emph{\myindex{first order ordinary differential equation}}
\index{ordinary differential equation}
\begin{equation*}
y' = F(x,y) .
\end{equation*}
Generally, we also specify an \emph{\myindex{initial condition}}
$y(x_0)=y_0$.  The solution of
the equation is a function $y(x)$ such that 
$y(x_0)=y_0$ and $y'(x) = F\bigl(x,y(x)\bigr)$.  See
\figureref{fig:sampleslopes} for a graphical representation of
this equation called a
\emph{\myindex{slope field}}.
\begin{myfigureht}
\myincludegraphics{sampleslopes}{%
A piece of the xy-plane near the origin with
a rectangular grid of slopes.  On the left half of the plot
the slopes go downward, but become flatter as we go up.
On the right half of the plot they point upward, and
again become flatter as we go up.
A point (x sub 0, y sub 0) is given in the first quadrant.
A graph of a solution is drawn that starts above
the x-axis goes down and then goes back up and through the
given point.}
\caption{A \emph{slope field} giving the slope $F(x,y)$ at each point,
in this case $F(x,y)=x(1-y)$.  A solution is drawn going through the point
$(x_0,y_0)=(1,0.3)$, notice how it follows the slopes.\label{fig:sampleslopes}}
\end{myfigureht}

When $F$ involves only the $x$ variable, the solution is given by the
fundamental theorem of calculus.  On the other hand, when $F$ depends
on both $x$ and $y$, we need far more firepower.  It is not always
true that a solution exists, and if it does, that it is the unique solution.
Picard's theorem gives us certain sufficient conditions for existence
and uniqueness.

\subsection{The theorem}

We need a definition of continuity in two variables.  A point in the
plane $\R^2 = \R \times \R$ is denoted by an ordered pair $(x,y)$.
For simplicity,
we give the following sequential definition of continuity.

\begin{defn}
Let $U \subset \R^2$ be a set, $F \colon U \to \R$ a function,
and $(x,y) \in U$ a point.  The function $F$ is
\emph{continuous at $(x,y)$}\index{continuous function of two variables}
if for every sequence
$\bigl\{ (x_n,y_n) \bigr\}_{n=1}^\infty$ of points in $U$ such that
$\lim_{n\to\infty} x_n = x$ and 
$\lim_{n\to\infty} y_n = y$, we have 
\begin{equation*}
\lim_{n \to \infty} F(x_n,y_n) = 
F(x,y) .
\end{equation*}
We say $F$ is \emph{continuous} if it is continuous at all points in $U$.
\end{defn}

\begin{thm}[Picard's theorem on existence and uniqueness]%
\index{existence and uniqueness theorem}\index{Picard's theorem}%
\label{thm:fs:picard}
Let $I, J \subset \R$ be closed bounded intervals, 
let $I^\circ$ and $J^\circ$ be their interiors%
\footnote{By interior of $[a,b]$, we mean $(a,b)$.},
and
let $(x_0,y_0) \in I^\circ \times J^\circ$.
Suppose $F \colon I \times J \to \R$ is continuous
and Lipschitz in the second variable, that is, there exists
an $L \in \R$ such that
\begin{equation*}
\babs{F(x,y) - F(x,z)} \leq L \sabs{y-z}
\qquad \text{for all } y,z \in J, x \in I .
\end{equation*}
Then there exists an $h > 0$ such that $[x_0-h,x_0+h] \subset I$ and a unique differentiable
function $f \colon [x_0 - h, x_0 + h] \to J \subset \R$ such that
\begin{equation} \label{picard:diffeq}
f'(x) = F\bigl(x,f(x)\bigr) \qquad \text{and} \qquad f(x_0) = y_0.
\end{equation}
\end{thm}

\begin{proof}
\pagebreak[2]
Suppose we could find a solution $f$.  Using the fundamental
theorem of calculus, we integrate the equation 
$f'(x) = F\bigl(x,f(x)\bigr)$, $f(x_0) = y_0$, and write \eqref{picard:diffeq}
as the integral equation
\begin{equation} \label{picard:inteq}
f(x) = y_0 + \int_{x_0}^x F\bigl(t,f(t)\bigr)\,dt .
\end{equation}
The idea of our proof is that we try to plug in approximations to
a solution to the right-hand side of \eqref{picard:inteq} to get better approximations on the
left-hand side of  \eqref{picard:inteq}.  We hope that in the end the sequence 
converges and solves
\eqref{picard:inteq} and hence \eqref{picard:diffeq}.
The technique below is called \emph{\myindex{Picard iteration}},
and the individual functions $f_k$ are called the 
\emph{Picard iterates}\index{Picard iterate}.

Without loss of generality, suppose $x_0 = 0$ (exercise below).
Another
exercise tells us that $F$ is bounded as it is continuous.
Therefore, pick some $M > 0$ so that 
$\babs{F(x,y)} \leq M$ for all $(x,y) \in I\times J$.
Pick $\alpha > 0$ such that
$[-\alpha,\alpha] \subset I$ and $[y_0-\alpha, y_0 + \alpha] \subset J$.
Define
\begin{equation*}
h \coloneqq \min \left\{ \alpha, \frac{\alpha}{M+L\alpha} \right\} .
\end{equation*}
Observe
$[-h,h] \subset I$.

Set $f_0(x) \coloneqq y_0$.
We define $f_k$ inductively.
Assuming $f_{k-1}([-h,h]) \subset [y_0-\alpha,y_0+\alpha]$,
we see 
$F\bigl(t,f_{k-1}(t)\bigr)$ is
a well-defined function of $t$ for $t \in [-h,h]$.
If $f_{k-1}$ is continuous
on $[-h,h]$, then
$F\bigl(t,f_{k-1}(t)\bigr)$ is
continuous as
a function of $t$ on $[-h,h]$ (left as an exercise).
Define
\begin{equation*}
f_k(x) \coloneqq y_0+ \int_{0}^x F\bigl(t,f_{k-1}(t)\bigr)\,dt ,
\end{equation*}
and $f_k$ is continuous on $[-h,h]$ by the fundamental theorem of calculus.
To see that $f_k$ maps $[-h,h]$ to $[y_0-\alpha,y_0+\alpha]$, we compute for
$x \in [-h,h]$
\begin{equation*}
\sabs{f_k(x) - y_0} = 
\abs{\int_{0}^x F\bigl(t,f_{k-1}(t)\bigr)\,dt }
\leq
M\sabs{x}
\leq
Mh
\leq
M
\frac{\alpha}{M+L\alpha}
\leq \alpha .
\end{equation*}
We next define $f_{k+1}$ using $f_k$ and so on.
Thus we have inductively defined a sequence $\{ f_k \}_{k=1}^\infty$
of continuous functions.
We need to show that it converges to a function $f$ that solves
the equation~\eqref{picard:inteq} and therefore~\eqref{picard:diffeq}.

We wish to show that the sequence $\{ f_k \}_{k=1}^\infty$ converges uniformly
to some function on $[-h,h]$.  First, for $t \in [-h,h]$,
we have the following
useful bound
\begin{equation*}
\Babs{F\bigl(t,f_{n}(t)\bigr) - 
F\bigl(t,f_{k}(t)\bigr)}
\leq
L \babs{f_n(t)-f_k(t)}
\leq
L \snorm{f_n-f_k}_{[-h,h]} ,
\end{equation*}
where $\snorm{f_n-f_k}_{[-h,h]}$ is the uniform norm,
that is, the supremum of $\babs{f_n(t)-f_k(t)}$ for $t \in [-h,h]$.
Now note that $\sabs{x} \leq h \leq \frac{\alpha}{M+L\alpha}$.
Therefore,
\begin{equation*}
\begin{split}
\babs{f_n(x) - f_k(x)}
& =
\abs{\int_{0}^x F\bigl(t,f_{n-1}(t)\bigr)\,dt 
-
\int_{0}^x F\bigl(t,f_{k-1}(t)\bigr)\,dt}
\\
& =
\abs{\int_{0}^x
\Bigl(
F\bigl(t,f_{n-1}(t)\bigr)-
F\bigl(t,f_{k-1}(t)\bigr)
\Bigr)
\,dt}
\\
& \leq
L
\,
\snorm{f_{n-1}-f_{k-1}}_{[-h,h]}
\,
\sabs{x}
\\
& \leq
\frac{L\alpha}{M+L\alpha}
\snorm{f_{n-1}-f_{k-1}}_{[-h,h]} .
\end{split}
\end{equation*}
Let $C \coloneqq \frac{L\alpha}{M+L\alpha}$ and note that $C < 1$.
Taking supremum on the left-hand side we get
\begin{equation*}
\snorm{f_n-f_k}_{[-h,h]} \leq C \snorm{f_{n-1}-f_{k-1}}_{[-h,h]} .
\end{equation*}
Without loss of generality,
suppose $n \geq k$.  Then by \hyperref[induction:thm]{induction} we can show 
\begin{equation*}
\snorm{f_n-f_k}_{[-h,h]} \leq C^{k} \snorm{f_{n-k}-f_{0}}_{[-h,h]} .
\end{equation*}
For $x \in [-h,h]$, we have
\begin{equation*}
\babs{f_{n-k}(x)-f_{0}(x)}
=
\babs{f_{n-k}(x)-y_0}
\leq \alpha .
\end{equation*}
Therefore,
\begin{equation*}
\snorm{f_n-f_k}_{[-h,h]} \leq C^{k} \snorm{f_{n-k}-f_{0}}_{[-h,h]} \leq C^{k} \alpha .
\end{equation*}
As $C < 1$, $\{f_n\}_{n=1}^\infty$ is uniformly Cauchy.
By \propref{prop:uniformcauchy}, $\{ f_n \}_{n=1}^\infty$
converges uniformly on $[-h,h]$ to some function $f \colon [-h,h] \to \R$.
The function $f$ is the uniform limit of continuous functions and thus
continuous.  Furthermore, since $f_n\bigl([-h,h]\bigr) \subset
[y_0-\alpha,y_0+\alpha]$ for all $n$,
then $f\bigl([-h,h]\bigr) \subset [y_0-\alpha,y_0+\alpha]$
(why?).


We now need to show that $f$ solves \eqref{picard:inteq}.
First, as before, we notice
\begin{equation*}
\Babs{F\bigl(t,f_{n}(t)\bigr) - 
F\bigl(t,f(t)\bigr)}
\leq
L \babs{f_n(t)-f(t)}
\leq
L \snorm{f_n-f}_{[-h,h]} .
\end{equation*}
As $\snorm{f_n-f}_{[-h,h]}$ converges to zero,
$F\bigl(t,f_n(t)\bigr)$ converges uniformly to $F\bigl(t,f(t)\bigr)$
for $t \in [-h,h]$.  Hence, for $x \in [-h,h]$
the convergence is uniform %(why?)
for $t \in [0,x]$ (or $[x,0]$ if $x < 0$).  Therefore,
\begin{align*}
y_0
+
\int_0^{x}
F(t,f(t)\bigr)\,dt
& =
y_0
+
\int_0^{x}
F\bigl(t,\lim_{n\to\infty} f_n(t)\bigr)\,dt
& &
\\
& =
y_0
+
\int_0^{x}
\lim_{n\to\infty} F\bigl(t,f_n(t)\bigr)\,dt
& & \text{(by continuity of } F\text{)}
\\
& =
\lim_{n\to\infty} 
\left(
y_0
+
\int_0^{x}
F\bigl(t,f_n(t)\bigr)\,dt
\right)
& & \text{(by uniform convergence)}
\\
& =
\lim_{n\to\infty} 
f_{n+1}(x)
=
f(x) .
& &
\end{align*}
We apply the fundamental theorem of calculus (\thmref{thm:FTCv2}) to show that
$f$ is differentiable and its derivative is $F\bigl(x,f(x)\bigr)$.  It is obvious
that $f(0) = y_0$.

Finally, what is left to do is to show uniqueness.  Suppose $g \colon [-h,h]
\to J \subset \R$ is another solution.
As before we use the fact that
$\babs{F\bigl(t,f(t)\bigr) - F\bigl(t,g(t)\bigr)} \leq L \snorm{f-g}_{[-h,h]}$.
Then
\begin{equation*}
\begin{split}
\babs{f(x)-g(x)}
& =
\abs{
y_0
+
\int_0^{x}
F\bigl(t,f(t)\bigr)\,dt
-
\left(
y_0
+
\int_0^{x}
F\bigl(t,g(t)\bigr)\,dt
\right)
}
\\
& =
\abs{
\int_0^{x}
\Bigl(
F\bigl(t,f(t)\bigr)
-
F\bigl(t,g(t)\bigr)
\Bigr)
\,dt
}
\\
& \leq
L\snorm{f-g}_{[-h,h]}\sabs{x}
\leq
Lh\snorm{f-g}_{[-h,h]}
\leq
\frac{L\alpha}{M+L\alpha}\snorm{f-g}_{[-h,h]} .
\end{split}
\end{equation*}
As 
before, $C = \frac{L\alpha}{M+L\alpha} < 1$.  By taking supremum over $x \in
[-h,h]$ on the
left-hand side we obtain
\begin{equation*}
\snorm{f-g}_{[-h,h]} \leq C \snorm{f-g}_{[-h,h]} .
\end{equation*}
This is only possible if $\snorm{f-g}_{[-h,h]} = 0$.  Therefore, $f=g$, and the
solution is unique.
\end{proof}

\subsection{Examples}

Let us look at some examples.  The proof of the theorem 
gives us an explicit way to find an $h$ that works.  It does not, however, give
us the best $h$.  It is often possible to find a much larger $h$ for
which the conclusion of the theorem holds.

The proof also gives us the Picard iterates as approximations to the
solution.  So the proof actually tells us how to obtain
the solution, not just that the solution exists.

\begin{example} \label{example:picardexponential}
Consider
\begin{equation*}
f'(x) = f(x), \qquad f(0) = 1 .
\end{equation*}
That is, we suppose $F(x,y) = y$, and we are looking for a function 
$f$ such that $f'(x) = f(x)$.
Let us forget for the moment that we solved this equation in 
\sectionref{sec:logandexp}.  See also \figureref{fig:exp} for a plot of
both the equation, showing
the slope $F(x,y)=y$ at each point, and the solution,
the exponential, that satisfies $f(0)=1$.

We pick any $I$ that contains 0
in the interior.
We pick an arbitrary $J$ that contains 1 in its interior.  We can
use $L = 1$.
The theorem guarantees an $h > 0$ such that
there exists a unique solution $f \colon [-h,h] \to \R$.  This solution
is usually denoted by
\begin{equation*}
e^x \coloneqq f(x) .
\end{equation*}
We leave it to the reader to verify that by picking $I$ and $J$
large enough the proof of the theorem guarantees that
we are able to pick $\alpha$ such that we get any
$h$ we want as long as $h < \nicefrac{1}{2}$.  We omit the calculation.
Of course, we know %(though we have not proved)
this function exists
as a function for all $x$, so an arbitrary $h$ ought to work, but the
theorem's proof only provides $h < \nicefrac{1}{2}$.

By the same reasoning as above,
no matter what $x_0$ and $y_0$ are,
the proof guarantees an arbitrary $h$ as long as $h < \nicefrac{1}{2}$.
Fix such an $h$.
We get a unique function defined on $[x_0-h,x_0+h]$.  After defining the
function on $[-h,h]$ we find a solution on the interval $[0,2h]$
and notice that the two functions must coincide on $[0,h]$ by uniqueness.
We thus iteratively construct the exponential for all $x \in \R$.
Therefore, Picard's theorem could be used to prove the existence and uniqueness
of the exponential.

Let us compute the Picard iterates.
We start with the constant function $f_0(x) \coloneqq 1$.  Then
\begin{align*}
f_1(x) & = 1 + \int_0^x f_0(s)\,ds =
1+x, \\
f_2(x) & = 1 + \int_0^x f_1(s)\,ds =
1 + \int_0^x (1+s)\,ds = 1 + x + \frac{x^2}{2}, \\
f_3(x) & = 1 + \int_0^x f_2(s)\,ds =
1 + \int_0^x \left(1+ s + \frac{s^2}{2} \right)\,ds =
1 + x + \frac{x^2}{2} + \frac{x^3}{6} .
\end{align*}
We recognize the beginning of the Taylor series for the exponential.
See \figureref{fig:exppicard}.
\begin{myfigureht}
\myincludegraphics{exppicardfig}{%
A graph of the exponential function as
a solid curve in the
region for x between a bit less than minus 1 and
a bit more than 1.
Several dashed curves give approximations.
They all go through the point (0,1).
First a horizontal line.  Then a line
with the same slope as the exponential at (0,1).
Then two further curved approximations that get yet closer
to the exponential, especially near x equals 0.}
\caption{The exponential (solid line) together with $f_0$, $f_1$, $f_2$,
$f_3$ (dashed).\label{fig:exppicard}}
\end{myfigureht}
\end{example}

\begin{example}
Consider the equation
\begin{equation*}
f'(x) = {\bigl(f(x)\bigr)}^2 \qquad \text{and} \qquad f(0)=1.
\end{equation*}
From elementary differential equations we know 
\begin{equation*}
f(x) = \frac{1}{1-x}
\end{equation*}
is the solution.
The solution is only defined on $(-\infty,1)$.  That is,
we are able to use $h < 1$, but never a larger $h$.
The function that takes $y$ to $y^2$ is
not Lipschitz as a function on all of $\R$.
As we approach $x=1$ from the left, the solution becomes larger
and larger.  The derivative of the solution grows as $y^2$, and so
the $L$ required has to be larger and larger as $y_0$ grows.
If we apply the
theorem with $x_0$ close to 1 and $y_0 = \frac{1}{1-x_0}$ we find
that the $h$ that the proof guarantees is smaller and smaller as $x_0$
approaches~1.

The $h$ from the proof is not the best $h$.
By picking $\alpha$ correctly, the proof of the theorem guarantees
$h=1-\nicefrac{\sqrt{3}}{2} \approx 0.134$ (we omit the
calculation) for $x_0=0$ and $y_0=1$, even though
we saw above that any $h < 1$ should work.
\end{example}

\begin{example}
Consider the equation
\begin{equation*}
f'(x) = 2 \sqrt{\babs{f(x)}}, \qquad f(0) = 0 .
\end{equation*}
The function $F(x,y) = 2 \sqrt{\sabs{y}}$ is continuous,
but not Lipschitz in $y$ (why?). 
The equation does not satisfy the hypotheses of the theorem.
The function
\begin{equation*}
f(x) =
\begin{cases}
x^2 & \text{if } x \geq 0,\\
-x^2 & \text{if } x < 0,
\end{cases}
\end{equation*}
is a solution, but $f(x) = 0$ is also a solution.
A solution exists, but is not unique.
\end{example}

\begin{example}
Consider $y' = \varphi(x)$ where $\varphi(x) \coloneqq 0$ if $x \in \Q$ and
$\varphi(x)\coloneqq 1$ if $x
\notin \Q$.  In other words, the $F(x,y) = \varphi(x)$ is discontinuous.
The equation has no solution regardless of the initial
conditions.
A solution would have
derivative $\varphi$, but $\varphi$ does not have the intermediate value property
at any point (why?).  No solution exists by
\hyperref[thm:darboux]{Darboux's theorem}.
\end{example}

The examples show that without the Lipschitz condition, a solution might
exist but not be unique, and without continuity of $F$, we may not
have a solution at all.
It is in fact a theorem, the Peano existence theorem, that if $F$ is
continuous a solution exists (but may not be unique).

\begin{remark}
It is possible to weaken what we mean by \myquote{solution to $y' = F(x,y)$}
by focusing on the integral equation
$f(x) = y_0 + \int_{x_0}^x F\bigl(t,f(t)\bigr) \, dt$.  For example,
let $H$ be the
Heaviside function%
\footnote{Named
for the English engineer, mathematician, and physicist
\href{https://en.wikipedia.org/wiki/Oliver_Heaviside}{Oliver Heaviside}
(1850--1925).},
that is $H(x) \coloneqq 0$ for $x < 0$ and $H(x) \coloneqq 1$ for $x \geq 0$.
Then
$y' = H(x)$, $y(0) = 0$,
is a common equation.
The \myquote{solution}
is the ramp function $f(x) \coloneqq 0$ if $x < 0$ and $f(x) \coloneqq x$ if $x \geq 0$,
since this function satisfies
$f(x) = \int_0^x H(t)\, dt$.  Notice, however, that $f'(0)$ does not exist,
so $f$ is only what is called a \emph{\myindex{weak solution}} to the
differential equation.
\end{remark}


\subsection{Exercises}

\begin{exercise}
Let $I, J \subset \R$ be intervals.
Let $F \colon I \times J \to \R$ be a continuous function
of two variables
and suppose $f \colon I \to J$ is a continuous function.
Show that $F\bigl(x,f(x)\bigr)$ is a continuous function on $I$.
\end{exercise}

\begin{exercise}
Let $I, J \subset \R$ be closed bounded intervals.
Show that if $F \colon I \times J \to \R$ is continuous,
then $F$ is bounded.
\end{exercise}

\begin{exercise}
We proved Picard's theorem under the assumption that $x_0 = 0$.
Prove the full statement of Picard's theorem for an arbitrary $x_0$.
\end{exercise}

\begin{exercise}
Let $f'(x)=x f(x)$ be our equation.  Start with the initial condition
$f(0)=2$ and find the Picard iterates $f_0,f_1,f_2,f_3,f_4$.
\end{exercise}

\begin{exercise}
Suppose $F \colon I \times J \to \R$
is a function that is continuous in the first variable,
that is, for every fixed $y$ the function that takes $x$ to $F(x,y)$ is
continuous.  Further, suppose $F$ is Lipschitz in the second variable,
that is, there exists a number $L$ such that
\begin{equation*}
\babs{F(x,y) - F(x,z)} \leq L \sabs{y-z}
\qquad \text{for all } y,z \in J, x \in I .
\end{equation*}
Show that $F$ is continuous as a function of two variables.  Therefore, the
hypotheses in the theorem could be made even weaker.
\end{exercise}

\begin{exercise}
\pagebreak[2]
A common type of equation is a
\emph{\myindex{linear first order differential equation}}, that is,
an equation of the form
\begin{equation*}
y' + p(x) y = q(x) , \qquad y(x_0) = y_0 .
\end{equation*}
Prove Picard's theorem for linear equations.  Suppose $I$ is an
interval, $x_0 \in I$, and $p \colon I \to \R$ and $q \colon I \to \R$ are
continuous.
Show that there exists a unique differentiable $f \colon I \to \R$,
such that $y = f(x)$
satisfies the equation and the initial condition.
Hint: Assume existence of the exponential function and use the integrating
factor formula for existence of $f$ (prove that it works and then that it is unique):
\begin{equation*}
f(x) \coloneqq e^{-\int_{x_0}^x p(s)\, ds} \left( \int_{x_0}^x e^{\int_{x_0}^t p(s)\, ds}
q(t) \,dt + y_0 \right).
\end{equation*}
\end{exercise}

\begin{exercise}
Consider the equation $f'(x) = f(x)$,
from \exampleref{example:picardexponential}.  Show that given any $x_0$,
any $y_0$, and any positive $h < \nicefrac{1}{2}$, we can pick $\alpha >
0$ large enough that the proof of Picard's theorem guarantees a solution
for the initial condition $f(x_0) = y_0$ in the interval
$[x_0-h,x_0+h]$.
\end{exercise}

\begin{exercise}
\pagebreak[3]
Consider the equation $y' = y^{1/3}x$.
\begin{enumerate}[a)]
\item
Show that for the initial condition $y(1)=1$, Picard's theorem applies.
Find an $\alpha > 0$, $M$, $L$, and $h$ that would work in the proof.
\item
Show that for the initial condition $y(1) = 0$, Picard's theorem
does not apply.
\item
Find a solution for $y(1) = 0$ anyway.
\end{enumerate}
\end{exercise}

\begin{exercise}
Consider the equation $x y' = 2y$.
\begin{enumerate}[a)]
\item
Show that $y = Cx^2$ is a solution for every constant $C$.
\item
Show that for every $x_0 \neq 0$ and every $y_0$, Picard's
theorem applies for the initial condition $y(x_0) = y_0$.
\item
Show that $y(0) = y_0$ is solvable if and only if $y_0 = 0$.
\end{enumerate}
\end{exercise}
